% !TEX program = pdflatex
\documentclass[11pt,a4paper]{article}

% --- Page margins and paragraphs ---
\usepackage[margin=1in]{geometry}
\usepackage{parskip}

% --- Basic math and typesetting ---
\usepackage{amsmath,amssymb,amsthm}
\usepackage{xcolor} %

% --- Theorem and definition environments ---
\theoremstyle{definition}
\newtheorem{definition}{Definition}
\newtheorem{example}{Example}

\theoremstyle{plain}
\newtheorem{theorem}{Theorem}
\newtheorem{proposition}{Proposition}
\newtheorem{lemma}{Lemma}

\usepackage{mathtools,booktabs,enumitem,graphicx,url}
\usepackage[hidelinks]{hyperref}
\setlength{\emergencystretch}{2em}
\setlist{topsep=4pt,itemsep=3pt,parsep=0pt}
\raggedbottom
\newcommand{\dd}{\,\mathrm{d}}
\newcommand{\C}{\mathbb{C}}

\begin{document}

% ==================================================
% Title and author information
% ==================================================
\begin{center}
    {\LARGE \textbf{Mathematics Method for Physics--Lecture 9}} \\[0.6em]
    {\large Laurent Series and Singularities: Exercises}
\end{center}

\vspace{-0.5em}
\noindent\rule{\textwidth}{0.5pt}
\vspace{0.3em}

\begin{center}
    Instructor: Zhaolong Han, SOCE, Shanghai Jiao Tong University;
\end{center}

\vspace{-0.6em}
\noindent Email: han.arkey@sjtu.edu.cn

\vspace{0.4em}
\noindent\rule{\textwidth}{0.5pt}

\vspace{0.8em}

These worked exercises accompany the notes on Laurent series and
singularities. Every expansion is stated together with its center
and convergence region. The logarithm example also specifies a branch.

\section{Taylor Expansions}
\subsection{Direct differentiation and indirect expansion}
For $f(z)=1/z^2$ about $z_0=1$, either differentiate $f$ repeatedly or use
\begin{equation}\label{eq:01}
\frac1{z^2}=-\left(\frac1z\right)'
 =-\left[\sum_{n=0}^{\infty}(-1)^n(z-1)^n\right]'.
\end{equation}
Reindexing gives
\begin{equation}\label{eq:02}
\boxed{\frac1{z^2}
 =\sum_{n=0}^{\infty}(-1)^n(n+1)(z-1)^n,\quad |z-1|<1.}
\end{equation}
The pole at $0$ fixes the radius $1$.

For $1/[z^2(z+1)]$, the poles are $0,-1$, so the radius about $1$ is
again $1$. For example, writing $w=z-1$ gives
\begin{equation}\label{eq:03}
\frac1{z^2(z+1)}
 =\left[\sum_{k=0}^{\infty}(-1)^k(k+1)w^k\right]
  \left[\frac12\sum_{j=0}^{\infty}\left(-\frac w2\right)^j\right],
 \qquad |w|<1.
\end{equation}
For $1/[z^2(z^2+1)]$, the poles are $0,\pm \mathrm{i}$; their distances from
$1$ are $1,\sqrt2,\sqrt2$, respectively. Its Taylor radius about $1$
is also $1$.

\section{Review: Laurent Expansions in a Specified Region}
\subsection{An exponential in the exterior of the unit circle}
Find the Laurent expansion about $0$ of
\begin{equation}\label{eq:04}
f(z)=\exp\left(\frac1{1-z}\right),\qquad |z|>1.
\end{equation}
Although
\begin{equation}\label{eq:05}
f(z)=\sum_{k=0}^{\infty}\frac1{k!}\left(\frac1{1-z}\right)^k,
\end{equation}
this is not yet a Laurent series in powers of $z$. For the specified
region, begin with
\begin{equation}\label{eq:06}
\frac1{1-z}=-\sum_{m=1}^{\infty}z^{-m},\qquad |z|>1.
\end{equation}
One can expand the powers and collect coefficients. A more systematic
method sets $w=1/z$, so the exterior becomes $0<|w|<1$:
\begin{equation}\label{eq:07}
F(w)\coloneqq f(1/w)=\exp\left(\frac w{w-1}\right).
\end{equation}
The apparent puncture at $w=0$ is removable; define $F(0)=1$. Then $F$
is holomorphic on $|w|<1$ and can be expanded by Taylor's theorem.

Write $F(w)=\sum_{n=0}^{\infty}a_nw^n$. Since
\begin{equation}\label{eq:08}
(1-w)^2 F'(w)+F(w)=0,
\end{equation}
coefficient comparison gives $a_0=1$, $a_1=-1$, and
\begin{equation}\label{eq:09}
(n+1)a_{n+1}+(1-2n)a_n+(n-1)a_{n-1}=0,\qquad n\geq1.
\end{equation}
Thus $a_2=-1/2$, $a_3=-1/6$, $a_4=1/24$, and
\begin{equation}\label{eq:10}
\boxed{e^{1/(1-z)}
   =1-\frac1z-\frac1{2z^2}-\frac1{6z^3}
        +\frac1{24z^4}+\cdots,\qquad |z|>1.}
\end{equation}
For a direct coefficient formula, expand
$F(w)=\exp[-w/(1-w)]$ to obtain, for $m\geq1$,
\begin{equation}\label{eq:11}
a_m=\sum_{k=1}^{m}\frac{(-1)^k}{k!}\binom{m-1}{k-1}.
\end{equation}

\subsection{A rational function on three regions}
Assume $|a|<|b|$ and let
\begin{equation}\label{eq:12}
g(z)=\frac1{(z-a)(z-b)}
     =\frac1{a-b}\left(\frac1{z-a}-\frac1{z-b}\right).
\end{equation}
The center and convergence region determine which geometric expansion
is appropriate.

\textbf{About $0$, for $|a|<|z|<|b|$:}
\begin{equation}\label{eq:13}
\boxed{g(z)=\frac1{a-b}\left[
       \sum_{n=0}^{\infty}a^nz^{-n-1}
       +\sum_{n=0}^{\infty}\frac{z^n}{b^{n+1}}\right].}
\end{equation}
Here $1/(z-a)$ is expanded in $a/z$, whereas $1/(z-b)$ is expanded in $z/b$.

\textbf{About $0$, in a neighborhood of infinity, $|z|>|b|$:}
\begin{equation}\label{eq:14}
\boxed{g(z)=\sum_{n=0}^{\infty}
               \frac{a^n-b^n}{a-b}\,z^{-n-1}.}
\end{equation}
The coefficient of $z^{-1}$ is zero because $a^0-b^0=0$.
The leading term is $z^{-2}$.

\textbf{About $a$, for $0<|z-a|<|a-b|$:} writing $u=z-a$ gives
\begin{equation}\label{eq:15}
\boxed{g(z)=\sum_{n=0}^{\infty}
             \frac{(-1)^n}{(a-b)^{n+1}}u^{n-1}.}
\end{equation}
This time the principal part is $1/[(a-b)(z-a)]$.

\section{Classification Exercises}
\subsection{Double poles of a rational function: 12(2)}
For
\begin{equation}\label{eq:16}
f(z)=\frac1{(z^2+\mathrm{i})^2},
\end{equation}
the roots of $z^2=-\mathrm{i}$ are
\begin{equation}\label{eq:17}
z_0=e^{-\mathrm{i}\pi/4},\qquad z_1=e^{3\mathrm{i}\pi/4}.
\end{equation}
Each is a simple zero of $z^2+\mathrm{i}$ and hence a double pole of $f$.
At infinity, $f(z)\sim z^{-4}$, so infinity is removable with value $0$.

\subsection{An essential singularity and simple poles: 12(6)}
Let
\begin{equation}\label{eq:18}
f(z)=\frac{e^{1/(z-1)}}{e^z-1}.
\end{equation}
At $z=2k\pi \mathrm{i}$, the denominator has a simple zero, and the numerator
is holomorphic and nonzero. Thus every such point, including $0$, is a
simple pole. At $1$, the denominator is nonzero and holomorphic.
The numerator has an essential singularity, so $1$ remains essential.
Infinity is nonisolated because the simple poles are unbounded.

\subsection{A sum with different types at \texorpdfstring{$\pm \mathrm{i}$}{+i and -i}: 12(10)}
For
\begin{equation}\label{eq:19}
f(z)=\frac1{1+z^2}+\cos\left(\frac1{z+\mathrm{i}}\right),
\end{equation}
$\mathrm{i}$ is a simple pole: the cosine term is holomorphic there.
At $-\mathrm{i}$, put $u=z+\mathrm{i}$. The cosine contributes
\begin{equation}\label{eq:20}
\cos(1/u)=\sum_{n=0}^{\infty}\frac{(-1)^n}{(2n)!}\,u^{-2n}.
\end{equation}
The rational term has only a simple pole at $-\mathrm{i}$, so it cannot cancel
infinitely many negative-power terms. Thus $-\mathrm{i}$ is essential.
At infinity the function has removable value $1$.

\subsection{A logarithm requires a branch choice: 12(12)}
Consider
\begin{equation}\label{eq:21}
f(z)=\frac1{z\log(1-z)}.
\end{equation}
Near $0$, take the holomorphic branch with $\log1=0$. Then
\begin{equation}\label{eq:22}
q(z)\coloneqq \frac1{f(z)}
      =z\log(1-z)=-z^2-\frac{z^3}{2}-\cdots .
\end{equation}
Hence $q$ has a zero of order $2$ at $0$, and $f$ has a double pole.
In particular,
\begin{equation}\label{eq:23}
f(z)=-\frac1{z^2}+\frac1{2z}+\frac1{12}+O(z).
\end{equation}
One may also check
\begin{equation}\label{eq:24}
q'(z)=\log(1-z)-\frac z{1-z},\qquad q''(0)=-2\neq0.
\end{equation}
At $1$, the logarithm has a branch point. There is no single-valued
holomorphic logarithm on a full punctured disk about $1$, so this point
does not fall under the removable/pole/essential classification.
The branch normalized at $0$ can, for example, be continued on
$\C\setminus[1,\infty)$.

\subsection{The extended plane: \texorpdfstring{$\cot z-1/z$}{cot(z)-1/z}}
Let
\begin{equation}\label{eq:25}
f(z)=\cot z-\frac1z
     =\frac{z\cos z-\sin z}{z\sin z}.
\end{equation}
At $0$, use Taylor expansions:
\begin{equation}\label{eq:26}
z\cos z-\sin z=-\frac{z^3}{3}+O(z^5),\qquad
 z\sin z=z^2+O(z^4).
\end{equation}
Thus
\begin{equation}\label{eq:27}
f(z)=-\frac z3+O(z^3)\longrightarrow0.
\end{equation}
The origin is removable, with extension $f(0)=0$.
At $k\pi$, $k\neq0$, $\sin z$ has a simple zero while $\cos(k\pi)$
is nonzero. Hence these points are simple poles of $\cot z$;
subtracting $1/z$, which is holomorphic there, does not remove them.
Their unbounded sequence makes infinity nonisolated.

\section{Assigned Exercises}
Textbook p.~117: 7(3), 7(4), 8, 11 and the assigned subitems of 12.

\section{Checks on the Calculations}
\begin{itemize}
\item Differentiating the geometric series gives
      $(1+w)^{-2}=\sum_{n\geq0}(-1)^n(n+1)w^n$ for $|w|<1$.
      The factor $n+1$ is needed in the expansion of $1/[z^2(z-1)]$.
\item The three-region example in the main notes uses the numerator $4$
      throughout. Its exterior expansion starts with $4(z-2)^{-2}$.
\item For $F(w)=\exp[w/(w-1)]$, $F^{(3)}(0)=-1$, so the coefficient of
      $z^{-3}$ in the exterior expansion is $-1/6$.
\item The derivative of $\log(1-z)$ is $-1/(1-z)$. On the branch with
      $\log1=0$, $1/[z\log(1-z)]$ has a double pole at $0$.
\item The expansion $\cot z-1/z=-z/3+O(z^3)$ gives the removable value
      $0$ at the origin.
\end{itemize}

\end{document}
